\documentclass[11pt,a4paper]{article}

% --- Packages ---
\usepackage[margin=.75in]{geometry}
\usepackage[utf8]{inputenc}
\usepackage[T1]{fontenc}
\usepackage{lmodern} % Fix for font not found error
\usepackage{titlesec}
\usepackage{enumitem}
\usepackage{hyperref}
\usepackage{xcolor}
\usepackage{fancyhdr}
\usepackage{comment}
\usepackage{cleveref}
\usepackage[bottom]{footmisc}
\usepackage{graphicx}
\usepackage{longtable}
\usepackage{booktabs}
\usepackage{array}
\usepackage{calc}
\usepackage{etoolbox}

% Standard Pandoc setup
\providecommand{\tightlist}{%
  \setlength{\itemsep}{0pt}\setlength{\parskip}{0pt}}

\crefformat{section}{\S#2#1#3}
\crefformat{subsection}{\S#2#1#3}
\crefformat{subsubsection}{\S#2#1#3}
\crefformat{figure}{#2Figure~#1#3}
\crefformat{footnote}{#2\footnotemark[#1]#3}

\linespread{0.9} % Riduce lo spazio tra le linee

% --- Colors ---
\definecolor{primary}{RGB}{0, 0, 0}
\definecolor{links}{HTML}{0000EE}
\definecolor{blue}{HTML}{26428B}

\hypersetup{
    colorlinks=true,
    linkcolor=links,
    urlcolor=links,
    citecolor=links,
}

% --- Section Formatting ---
\titleformat{\section}
  {\normalfont\Large\bfseries\color{blue}}{\thesection}{1em}{}
\titleformat{\subsection}
{\normalfont\large\bfseries\color{blue}}{\thesubsection}{1em}{}
\titleformat{\subsubsection}
{\normalfont\bfseries\color{blue}}{\thesubsubsection}{1em}{}

% --- Header and Footer ---
\pagestyle{fancy}
\fancyhead[L]{\small Giuseppe Capasso}
\fancyhead[R]{\small intro-to-linux}
\fancyfoot[C]{\thepage}
\renewcommand{\headrulewidth}{0.4pt}

% --- Custom Title Command ---
\newcommand{\proposaltitle}[1]{
    \begin{center}
        \vspace*{0.1cm}
        {\LARGE \bfseries \color{blue} #1} \\[1.5ex]
        {\large \textbf{Giuseppe Capasso}} \\[1ex]
                \small{
            \vspace{0.1cm}
            \textbf{Materia:} intro-to-linux\\
        }
            \end{center}
    \vspace{0.3cm}
    \hrule height 0.5pt
    \vspace{0.5cm}
}

\begin{document}

\proposaltitle{Test Finale: Introduzione a Linux}

Test Finale: Introduzione a Linux Nome e Cognome:
\_\_\_\_\_\_\_\_\_\_\_\_\_\_\_\_\_\_\_\_\_\_\_\_\_\_\_\_\_\_\_\_\_\_\_\_\_\_\_\_\_\_\_\_\_\_\_\_\_\_\_\_\_
Data: \_\_\_\_\_\_\_\_\_\_\_\_\_\_\_\_\_\_\_

Punteggio: \_\_\_\_\_/20

\begin{enumerate}
\def\labelenumi{\arabic{enumi}.}
\tightlist
\item
  Qual è la differenza tecnica tra ``Linux'' e una ``Distribuzione''
  (es. Ubuntu)? A. Linux è il sistema operativo completo, la
  distribuzione è solo l'interfaccia grafica. B. Linux è il Kernel (il
  cuore che gestisce l'hardware), la distribuzione è il pacchetto
  completo di software e strumenti. C. Sono sinonimi, indicano la stessa
  cosa. D. Linux è un software a pagamento, le distribuzioni sono sempre
  gratuite.
\item
  Quale comando utilizzi per visualizzare il percorso completo della
  directory in cui ti trovi (es. /home/student)? A. ls B. cd C. pwd D.
  whoami
\item
  Devi rinominare il file bozza.txt in finale.txt. Quale comando
  utilizzi? A. rn bozza.txt finale.txt B. cp bozza.txt finale.txt C. mv
  bozza.txt finale.txt D. rm bozza.txt finale.txt
\item
  In un ambiente di virtualizzazione, cos'è l'\,``Hypervisor''? A. Il
  sistema operativo installato dentro la macchina virtuale. B. Il
  software che crea e gestisce le macchine virtuali (es. VirtualBox). C.
  Un virus che colpisce solo i sistemi Linux. D. L'amministratore di
  sistema con privilegi elevati.
\item
  Cosa succede se esegui cd ..? A. Ti sposti nella directory Home
  dell'utente.
\end{enumerate}

B. Cancelli la directory corrente. C. Ti sposti nella directory di
livello superiore (cartella genitore). D. Il sistema ti chiede di
confermare l'uscita dal terminale. 6. Secondo lo standard FHS, quale
directory contiene i file di configurazione del sistema? A. /bin B.
/home C. /var D. /etc 7. Cosa esegue esattamente il comando sudo apt
update? A. Aggiorna tutti i programmi installati all'ultima versione. B.
Scarica l'elenco aggiornato dei pacchetti dai repository (aggiorna
l'indice). C. Aggiorna il Kernel di Linux. D. Rimuove i pacchetti non
più necessari. 8. Analizzando i permessi drwxr-xr--, quali diritti hanno
gli ``Altri'' (Others)? A. Lettura ed Esecuzione. B. Scrittura ed
Esecuzione. C. Sola Lettura. D. Nessun permesso. 9. Vuoi assegnare
permessi completi (rwx) al proprietario e solo lettura/esecuzione (r-x)
a tutti gli altri. Qual è il comando corretto (supponi il file abbia I
seguenti permessi ``r--r--r--''? A. chmod u+wx,o+x,g+x file B. chmod
u+x,o-wrx,g-x file C. chmod addx file D. chmod default file 10. Qual è
il file che contiene la lista degli utenti del sistema? A. /etc/shadow
B. /etc/passwd

C. /etc/users D. /home/list.txt 11. Come viene identificata la seconda
partizione del primo disco SATA in Linux? A. /dev/hda2 B. /dev/sdb1 C.
D:\\
D. /dev/sda2 12. Quale file devi modificare per rendere permanente il
mount di un disco al riavvio del sistema? A. /etc/mount B. /etc/fstab C.
/dev/disk D. /etc/modules 13. Prima di utilizzare una nuova partizione
per salvarci dei dati, quale operazione è obbligatoria? A. Copiarci
dentro un file di testo. B. Creare un File System (Formattazione), ad
esempio con mkfs.ext4. C. Riavviare la macchina virtuale. D. Installare
il driver del disco. 14. Cos'è un ``Backup Incrementale''? A. Una copia
completa di tutti i dati, fatta ogni giorno. B. Una copia solo dei file
modificati rispetto all'ultimo backup (di qualsiasi tipo). C. Una copia
che sovrascrive i dati originali. D. Una copia dei soli file di sistema.
15. Quale comando è ideale per sincronizzare due cartelle copiando solo
le differenze? A. cp -r B. tar -czvf C. rsync -av D. mv

16. Qual è il comando più moderno per visualizzare gli indirizzi IP
delle interfacce di rete? A. ipconfig B. ifconfig C. ip addr show (o ip
a) D. netstat 17. A quale livello opera un indirizzo MAC? A. Livello 3
(IP / Routing). B. Livello 2 (Fisico / Data Link -- Switch). C. Livello
4 (Porte / Trasporto). D. Livello 7 (Applicazione). 18. Vuoi associare
manualmente il nome server-test all'IP 192.168.1.50 senza usare un DNS
esterno. Quale file modifichi? A. /etc/resolv.conf B. /etc/hosts C.
/etc/network/interfaces D. /etc/hostname 19. Qual è la porta standard
per il traffico Web non criptato (HTTP)? A. 22 B.443 C. 80 D. 8080 20.
Qual è l'indirizzo IP di ``Loopback'' (che identifica sempre il computer
stesso)? A. 192.168.1.1 B. 127.0.0.1 C. 0.0.0.0 D. 10.0.0.1

Domanda Risposta Argomento Note Rapide 1

B

Intro

Kernel vs Distro

2

C

CLI

pwd = Print Working Directory

3

C

CLI

mv sposta o rinomina

4

B

Virt

Hypervisor gestisce le VM

5

C

CLI

.. = Directory genitore

6

D

Admin

/etc = Editable Text Config

7

B

APT

update aggiorna solo la lista

8

C

Permessi

r-- = Solo lettura (4)

9

A

Permessi

chmod u+wx,o+x,g+x

10

B

Utenti

passwd contiene gli utenti, shadow le psw

11

D

Storage

sda=Disco 1, 2=Partizione 2

12

B

Storage

fstab = Filesystem Table

13

B

Storage

Senza Filesystem non scrivi dati

14

B

Backup

Risparmia spazio e tempo

15

C

Backup

rsync è standard per sync differenziale

16

C

Network

ip sostituisce ifconfig

17

B

Network

MAC è indirizzo fisico (L2)

18

B

Network

/etc/hosts è il DNS locale

19

C

Network

80=HTTP, 443=HTTPS, 22=SSH

20

B

Network

Localhost standard



\end{document}
